\begin{deluxetable*}{lcccc}
\tablecaption{Post hoc: contribution of each test. \label{tab:abl}}
\tablewidth{0pt}
\tablehead{\colhead{} & \colhead{Simple, B1(i)} & \colhead{Simple, B5} & \colhead{Learned, B1(i)} & \colhead{Learned, B5}}
\startdata
All six terms & 0.775 & 0.800 & \nodata & \nodata \\
without F1, predicted event coverage & 0.316 ($-$0.459) & 0.241 ($-$0.559) & $-$0.015 & $+$0.000 \\
without F2, odd/even depth & 0.691 ($-$0.085) & 0.475 ($-$0.324) & $-$0.004 & $-$0.007 \\
without F3, secondary eclipse & 0.748 ($-$0.028) & 0.667 ($-$0.133) & $-$0.023 & $-$0.004 \\
without F4, stellar density & 0.772 ($-$0.003) & 0.818 ($+$0.019) & $-$0.108 & $+$0.009 \\
without F5, ephemeris coherence & 0.785 ($+$0.009) & 0.800 ($+$0.001) & $-$0.012 & $-$0.007 \\
without F6, shape of alternating events & 0.776 ($+$0.001) & 0.800 ($+$0.000) & $-$0.007 & $+$0.002 \\
simple posterior permuted & \nodata & \nodata & $-$0.525 & $-$0.298
\enddata
\tablecomments{Columns 2 to 3: accuracy of the simple combiner with one test left out (change in parentheses), on the B1(i) test injections and the B5 primary sample. Columns 4 to 5: change in accuracy of the frozen learned combiner when the test and its inputs are permuted between targets (mean of five permutations). The learned combiner also sees the simple posterior, which contains all six terms, so its single test values are lower limits. Nothing was refitted.}
\end{deluxetable*}
